% TOP_PROBLEM: {"id": 3166, "title": "Barnette's Conjecture", "category_id": 3, "category": {"id": 3, "name": "graph_theory", "display_name": "Graph Theory", "description": "Problems involving graphs, networks, and their properties.", "slug": "graph-theory", "order_index": 3, "created_at": "2026-07-31T15:26:25.671Z"}, "status": "open", "problem_number": "OPG-385", "set_id": 12}
% FIRST_POSTED: 2026-10-06
% ATTEMPT_STATUS: solved
% OPENAI_CLAIM: {"schema_version":1,"source_repo":"https://github.com/openai/math","source_commit":"adc7f1241b42e322a6451854ab7e4b4c146bf78a","release_date":"2026-10-06","match":"full","resolution":"proved","independently_reviewed":false,"second_pass":"agreed","families":[{"family":"180","title":"Barnette’s Hamiltonian-cycle conjecture","subject":"Combinatorics","manuscripts":[{"title":"Paired states and Hamiltonian cycles in cubic bipartite planar graphs","pdf_path":"preprints/Paired-states-and-Hamiltonian-cycles-in-cubic-bipartite-planar-graphs-September-24-2026/paper.pdf","readme_path":"preprints/Paired-states-and-Hamiltonian-cycles-in-cubic-bipartite-planar-graphs-September-24-2026/README.md","date":"2026-09-24","theorem_ref":"thm:main (paper.tex:79); thm:partition (paper.tex:208); cor:edge-avoidance (paper.tex:794); cor:pfaffian-paths (paper.tex:812)","statement_sources":[{"path":"preprints/Paired-states-and-Hamiltonian-cycles-in-cubic-bipartite-planar-graphs-September-24-2026/build/paper.tex","line":79,"label":"thm:main"},{"path":"preprints/Paired-states-and-Hamiltonian-cycles-in-cubic-bipartite-planar-graphs-September-24-2026/build/paper.tex","line":208,"label":"thm:partition"},{"path":"preprints/Paired-states-and-Hamiltonian-cycles-in-cubic-bipartite-planar-graphs-September-24-2026/build/paper.tex","line":794,"label":"cor:edge-avoidance"},{"path":"preprints/Paired-states-and-Hamiltonian-cycles-in-cubic-bipartite-planar-graphs-September-24-2026/build/paper.tex","line":812,"label":"cor:pfaffian-paths"}]}],"lean":{"doc_path":"lean/docs/180.md","comparators":["lean/ComparatorChallenges/BarnetteHamiltonian.json","lean/ComparatorChallenges/BarnetteHamiltonian.lean"],"declarations":[],"covers_main_theorem":true},"reasoning_trace":null}]}
\documentclass[11pt]{article}
\usepackage[margin=25mm]{geometry}
\usepackage{fontspec}
\setmainfont{Latin Modern Roman}
\usepackage{amsmath,amssymb,mathtools}
\usepackage{unicode-math}
\setmathfont{Latin Modern Math}
\usepackage{enumitem}
\usepackage{xurl,hyperref,bookmark}
\hypersetup{colorlinks=true,urlcolor=blue}
\setcounter{secnumdepth}{0}
\setlength{\parindent}{0pt}
\setlength{\parskip}{5pt}
\setlength{\emergencystretch}{3em}
\begin{document}
\section*{Barnette's Conjecture}

\subsection{OpenAI's claim}
\textit{Barnette Hamiltonian cycles}: The claimed Hamiltonicity theorem concerns finite simple cubic bipartite planar graphs with vertex connectivity three. It also yields avoidance of an individually prescribed edge by a suitable Hamiltonian cycle.

\subsection{Mathematical statement of the claimed result}
\paragraph{Mathematical claim 1.} For a finite graph without loops or multiple edges, assume planarity, bipartiteness, degree three at every vertex, and $3$-vertex-connectivity. The asserted conclusion is existence of a Hamiltonian cycle, that is, a cycle passing through every vertex exactly once. A stated strengthening says that for each chosen edge $e$ there is such a cycle avoiding $e$. The planar dual formulation partitions the vertices of a black-and-white face-colored sphere triangulation into two induced trees, with either specified facial-edge endpoints in one tree and the third face vertex in the other.

\subsection{Outline of the argument}
\textit{Barnette Hamiltonian cycles}: The sphere dual is a triangulation with black and white triangular faces. Hamiltonicity is translated into a partition of its vertices into two induced trees, where every within-class edge must be retained. Paired states organize local modifications and the prescribed facial-edge conditions, producing the partition and then the primal cycle.

\subsection{Scope relative to this problem}
\textit{Barnette Hamiltonian cycles}: The theorem has precisely the cubic, bipartite, planar and three-vertex-connected hypotheses of Barnette’s conjecture and yields a Hamiltonian cycle.

With the hypotheses and conventions above, the claimed result addresses the whole mathematical assertion on this page.

\subsection{Formal verification}
Family 180 describes Lean coverage for Hamiltonian-cycle existence in the cubic bipartite planar vertex-connected setting.

\textit{Barnette Hamiltonian cycles}: The manuscript links \texttt{BarnetteHamiltonian.lean}.

The documented selection includes this manuscript’s main assertion; that assertion retains the restrictions stated above.

This site has not run the Lean development. These links describe OpenAI’s selected formal statements and should be read with their stated scope.

\subsection{Status caveat}
OpenAI reports that an internal model produced the manuscripts. Their arguments have not been independently verified as proofs by this site. Any portion outside the documented formal scope may contain mistakes, and the catalogue label reports the announced mathematical claim.

\subsection{Sources}
\begin{itemize}
\item OpenAI, \emph{Paired states and Hamiltonian cycles in cubic bipartite planar graphs}, 2026-09-24; citation key \texttt{OAI:Paired-states-and-Hamiltonian-cycles-in-cubic-bipartite-planar-graphs-September-24-2026}. \href{https://github.com/openai/math/blob/adc7f1241b42e322a6451854ab7e4b4c146bf78a/preprints/Paired-states-and-Hamiltonian-cycles-in-cubic-bipartite-planar-graphs-September-24-2026/paper.pdf}{PDF (GitHub)}; \href{https://github.com/openai/math/raw/adc7f1241b42e322a6451854ab7e4b4c146bf78a/preprints/Paired-states-and-Hamiltonian-cycles-in-cubic-bipartite-planar-graphs-September-24-2026/paper.pdf}{PDF (direct)}; \href{https://github.com/openai/math/blob/adc7f1241b42e322a6451854ab7e4b4c146bf78a/preprints/Paired-states-and-Hamiltonian-cycles-in-cubic-bipartite-planar-graphs-September-24-2026/README.md}{Manuscript page and citation}; \href{https://github.com/openai/math/blob/main/preprints/Paired-states-and-Hamiltonian-cycles-in-cubic-bipartite-planar-graphs-September-24-2026/paper.pdf}{Latest version}; Statement source: \href{https://github.com/openai/math/blob/adc7f1241b42e322a6451854ab7e4b4c146bf78a/preprints/Paired-states-and-Hamiltonian-cycles-in-cubic-bipartite-planar-graphs-September-24-2026/build/paper.tex\#L79}{thm:main}; Statement source: \href{https://github.com/openai/math/blob/adc7f1241b42e322a6451854ab7e4b4c146bf78a/preprints/Paired-states-and-Hamiltonian-cycles-in-cubic-bipartite-planar-graphs-September-24-2026/build/paper.tex\#L208}{thm:partition}; Statement source: \href{https://github.com/openai/math/blob/adc7f1241b42e322a6451854ab7e4b4c146bf78a/preprints/Paired-states-and-Hamiltonian-cycles-in-cubic-bipartite-planar-graphs-September-24-2026/build/paper.tex\#L794}{cor:edge-avoidance}; Statement source: \href{https://github.com/openai/math/blob/adc7f1241b42e322a6451854ab7e4b4c146bf78a/preprints/Paired-states-and-Hamiltonian-cycles-in-cubic-bipartite-planar-graphs-September-24-2026/build/paper.tex\#L812}{cor:pfaffian-paths}.
\item Family 180: \href{https://github.com/openai/math/blob/adc7f1241b42e322a6451854ab7e4b4c146bf78a/lean/docs/180.md}{Lean scope}.
\item \href{https://github.com/openai/math/blob/adc7f1241b42e322a6451854ab7e4b4c146bf78a/lean/ComparatorChallenges/BarnetteHamiltonian.json}{Comparator BarnetteHamiltonian.json}.
\item \href{https://github.com/openai/math/blob/adc7f1241b42e322a6451854ab7e4b4c146bf78a/lean/ComparatorChallenges/BarnetteHamiltonian.lean}{Comparator BarnetteHamiltonian.lean}.
\item \href{https://github.com/openai/math/blob/adc7f1241b42e322a6451854ab7e4b4c146bf78a/CONTENTS.md}{OpenAI catalogue entry}.
\end{itemize}
\end{document}
